\section{Data and Valuation Setup}

\subsection{Controlled layer separation}
The integrated experiment is organized as a controlled handoff across three modules. Team 8 governs only the gold-price process; Team 4 supplies the operating forecasts; and Team 5 maps those inputs through a common tax, WACC, terminal-value contract without a reconciled equity bridge. In the four-model valuation comparison, only the gold engine changes. Barrick's Q1--Q2 2026 actuals and the same Team 4 production and unit-cost forecasts from Q3 2026 onward are frozen over 20 quarters; the same Team 5 corporate rules and WACC shocks are then reused across runs.



\subsection{Calibration snapshot and measure discipline}
The current calibration uses the 2 September 2026 GLD call surface acquired through the IB API. After the common DTE $\geq75$ and numerical filters, the dense surface contains 605 eligible observations across 12 expiries and 146 distinct strikes, with 79--653 days to expiry. The official $8\times8$ Chebyshev--Chebyshev sample contains 64 distinct actual contracts spanning eight expiries and 20 strikes. The risk-free input is a same-date NSS fit to 14 U.S. Treasury tenors \citep{nss1994,treasury2026}. Its 2.0608 bp RMSE is computed against the continuously compounded par-yield proxy; the curve is not described as a bootstrapped zero curve. The Barrick market reference is the 2 September close of USD~44.13, while the common gold anchor remains Barrick's Q2 2026 average realized price of USD~4,417/oz \citep{barrick2026}. The operating calendar is a fixed scenario index, not a forecast commencing on 2 September; earlier actuals are not future observations.

The option layer is calibrated under $\mathbb Q$, with deterministic pricing rates. The operating experiment instead defines an artificial scenario law $\mathcal R$: calibrated coefficients are reanchored, commodity mean schedules imposed and independent WACC shocks sampled. Neither a Radon--Nikodym density nor physical risk premia are estimated. Consequently $\mathcal R$ is neither a validated $\mathbb P$ forecast nor the original traded-instrument $\mathbb Q$ law. Gold uses GLD shape at a separately dated realised-gold level; copper uses a cross-delivery mean schedule. Neither transformation is a change of measure.

For a fixed-delivery futures price, the deterministic-rate pricing convention has zero drift. Copper's delivery-curve slope is a scenario assumption, not that futures time drift. Jump compensation uses the simulated intensity and mark moment $\lambda(\E[e^J]-1)$. A genuine $\mathbb Q\rightarrow\mathbb P$ specification would need price and variance risk premia and, with jumps, intensity and mark-law adjustments; changing drift alone would not estimate them. Futures are also not generally physical expected spot prices \citep{schwartz1997}.

Ordinary DCF uses physical expected cash flows with a compatible risk-adjusted discount rate; a pricing-measure valuation instead needs its consistent numeraire and pricing kernel. Q-shaped scenarios plus assumed WACC may mix incompatible risk adjustments. Replacing WACC by a risk-free curve would not by itself reconcile this incomplete corporate claim. WACC shocks are not short-rate shocks. The reported operator is therefore a conditional sensitivity, without economic valuation or probability interpretation.

\begin{table}[H]
\centering
\caption{Frozen contract used for model attribution.}
\label{tab:frozen_contract}
\begin{tabularx}{\columnwidth}{@{}>{\bfseries}p{0.25\columnwidth}X@{}}
\toprule
Block & Current experiment \\
\midrule
Gold layer & Black--Scholes/GBM, Heston, Bates--Poisson and Full Bates--Hawkes calibrated by Team 8 under $\mathbb Q$. \\
Operations & Barrick Q1--Q2 2026 actuals followed by Team 4 deterministic production and unit-cost vectors from Q3 2026. \\
Corporate mapping & Team 5 tax, growth--ROIC--reinvestment schedule, stochastic WACC and explicit terminal convention; no equity bridge. \\
Simulation & 8,192 paths, 260 fine steps, 20 quarters and a common WACC shock array. \\
\bottomrule
\end{tabularx}
\end{table}

\subsection{Option exercise convention}
GLD contracts are American-style \citep{cboe2026}. We set payout yield $q=0$: trust expenses reduce backing, not cash distributions \citep{spdr2025}. For non-distributing calls and positive rates, early exercise is suboptimal in the frictionless model. An 800-step CRR check on all 64 contracts gives an American--European premium below $10^{-8}$ USD; this tests that convention, not market frictions.

\subsection{Operating-model status}
Team 4's original operating study is richer than the object ultimately passed downstream. Cost of Sales is forecast from a log-space master chain; production is constructed bottom-up from ore processed, average grade and recovery. The unified valuation refreshes the handoff with 2026 actuals but still freezes the resulting 20-quarter production and cost vectors. Consequently, mine-level operating uncertainty is not propagated jointly with gold-price uncertainty. Team 4's legacy price simulation and illustrative EBITDA/valuation output are explicitly excluded.

\subsection{Gold perimeter and reporting basis}
The Team 4 subset comprises Bulyanhulu, Carlin, Cortez, Kibali, Loulo--Gounkoto, North Mara, Pueblo Viejo and Turquoise Ridge. Its quarterly estimation window is 2019Q1--2025Q4, with North Mara beginning in 2019Q4. The operating quantities are already attributable: Nevada Gold Mines is 61.5\%, Pueblo Viejo 60\%, Kibali 45\%, Loulo--Gounkoto 80\%, and the Tanzanian mines 84\%. Kibali is equity-accounted, whereas the other admitted interests are consolidated with partners' non-controlling interests \citep{barrickar2025}. Those partners must not be deducted again from attributable flows. Tanzania's contractual economic-benefit sharing still requires a separate cash-flow reconciliation.

The frozen handoff mixes company-wide gold actuals of 719 and 796 koz in Q1--Q2 2026 with later forecasts for the eight-mine subset; that subset produced only 622 and 711 koz in those quarters. The scope break is retained and disclosed, not silently described as a consistent mine panel. Cost of Sales includes depreciation and price-sensitive royalties; multiplying its sold-ounce cost by production yields a component-margin proxy, not EBITDA or reconciled EBIT. Loulo--Gounkoto's 2025Q4 accounting uplift also makes its cost anchor a normalization issue. Ownership, the North America IPO, cash/debt, share count and reserves must be frozen at a common corporate valuation date before a per-share result is permitted.
